\chapter{Comparison with Other Approaches}
\label{ch:comparison}

Veilbox occupies a unique niche in the Linux distribution landscape. This
chapter compares Veilbox with other approaches to building minimal Linux
systems for container hosting.

\section{Comparison Overview}

\begin{longtable}{|p{3cm}|p{3cm}|p{3cm}|p{3cm}|p{3cm}|}
\hline
\textbf{Feature} & \textbf{Veilbox} & \textbf{Alpine} & \textbf{Buildroot} & \textbf{Fedora CoreOS} \\
\hline
Kernel source & Custom & Mainline & Custom & Mainline \\
Init system & BusyBox init & OpenRC & BusyBox init / systemd & systemd \\
Package manager & None & apk & None & rpm-ostree \\
Rootfs size & ~90 MB & ~50 MB & ~20 MB & ~1 GB \\
Build time & ~5 min & N/A & ~15 min & N/A \\
Persistent storage & State partition & Full disk & Full disk & Full disk \\
Container runtime & containerd & Docker / containerd & Docker & podman \\
Config language & Bash \& Kconfig & n/a & Kconfig & Ignition \\
Minimum RAM & 2 GB & 128 MB & 64 MB & 1 GB \\
\hline
\end{longtable}

\section{Veilbox vs. Alpine Linux}

Alpine Linux is the most popular minimal container host OS:

\begin{itemize}[nosep]
  \item \textbf{Similarities}: Both use BusyBox for core utilities. Both
        aim for minimal footprint. Both target container workloads.
  - \textbf{Differences}: Alpine uses OpenRC init (more features), has
        the apk package manager, uses musl libc (smaller but not always
        compatible), and is a general-purpose distribution with hundreds
        of packages. Veilbox has no package manager, uses glibc, and is
        single-purpose.
  - \textbf{Trade-off}: Alpine is more flexible (install packages at
        runtime). Veilbox is more secure (no runtime installation
        possible) and more deterministic (every deployment is identical).
\end{itemize}

\section{Veilbox vs. Buildroot}

Buildroot is a build system for embedded Linux:

\begin{itemize}[nosep]
  \item \textbf{Similarities}: Both build a custom kernel with embedded
        rootfs. Both support BusyBox. Both produce minimal images.
  \item \textbf{Differences}: Buildroot has a sophisticated package selection
        system, supports hundreds of packages, and targets embedded
        hardware (ARM, MIPS, RISC-V). Veilbox is x86\_64-only with a
        fixed set of packages.
  \item \textbf{Trade-off}: Buildroot is more powerful but requires learning
        its configuration system and has a longer initial build.
        Veilbox is simpler to configure (just edit build.sh).
\end{itemize}

\section{Veilbox vs. Fedora CoreOS}

Fedora CoreOS is Red Hat's minimal container OS:

\begin{itemize}[nosep]
  \item \textbf{Similarities}: Both target container workloads. Both use
        containerd (or podman). Both support automated updates.
  \item \textbf{Differences}: CoreOS uses systemd (complex), Ignition for
        provisioning (YAML-based), rpm-ostree for atomic updates, and
        has SELinux enabled by default. It is supported by Red Hat.
        Veilbox is independently built, uses BusyBox init, and has no
        update mechanism.
  \item \textbf{Trade-off}: CoreOS is enterprise-ready with support and
        updates. Veilbox is simpler (no systemd, no SELinux) and easier
        to understand entirely.
\end{itemize}

\section{When to Use Veilbox}

Veilbox is ideal for:
\begin{itemize}[nosep]
  \item \textbf{Learning}: Understanding every component in a Linux system.
  \item \textbf{CI/CD}: Disposable VMs for testing containers.
  \item \textbf{Education}: Teaching kernel compilation, boot process, and
        init systems.
  \item \textbf{Benchmarking}: Minimal overhead for container performance
        testing.
\end{itemize}


\section{Detailed Feature Matrix}

\begin{longtable}{|p{4cm}|p{4cm}|p{4cm}|}
\hline
\textbf{Feature} & \textbf{Veilbox} & \textbf{Typical Distro} \\
\hline
Init system & BusyBox init & systemd / OpenRC \\
Process supervisor & None (busybox init) & systemd / supervise-daemon \\
Service files & Shell scripts & systemd unit files \\
Dependency-based boot & No & Yes \\
Parallel boot & No & Yes \\
Socket activation & No & Yes \\
Timer units & No & Yes \\
Logging & syslogd + /var/log/messages & journald \\
User session management & No & logind \\
Tmpfiles & No & systemd-tmpfiles \\
Sysctl & Manual (/proc/sys) & systemd-sysctl \\
Mount units & No & systemd-mount \\
Timer units & No & systemd-timer \\
\hline
\end{longtable}

\section{Size Comparison Breakdown}

\begin{longtable}{|p{4cm}|p{3cm}|p{3cm}|p{3cm}|}
\hline
\textbf{Component} & \textbf{Veilbox} & \textbf{Alpine} & \textbf{Ubuntu Server} \\
\hline
Kernel & 10 MB & 8 MB & 12 MB \\
Initramfs & 80 MB & 15 MB & 50 MB \\
Rootfs (installed) & N/A (embedded) & 50 MB & 2 GB \\
BusyBox & 1.2 MB & 1.2 MB & N/A \\
coreutils & N/A (BusyBox) & N/A (BusyBox) & 15 MB \\
systemd & N/A & N/A & 25 MB \\
OpenRC & N/A & 2 MB & N/A \\
Package manager & None & apk (3 MB) & apt (20 MB) \\
Container runtime & containerd (35 MB) & Docker (100 MB) & Docker (100 MB) \\
\hline
\textbf{Total} & \textbf{~90 MB} & \textbf{~180 MB} & \textbf{~2.5 GB} \\
\hline
\end{longtable}

\section{Use Case Recommendations}

\begin{itemize}[nosep]
  \item \textbf{Veilbox}: Learning, CI/CD VMs, container testing,
        security research, embedded container hosts.
  \item \textbf{Alpine Linux}: Production container hosts, Docker hosts,
        general-purpose embedded systems, routers, firewalls.
  \item \textbf{Buildroot}: Custom embedded devices, IoT, specialized
        hardware, minimal appliances.
  \item \textbf{Fedora CoreOS}: Enterprise container hosts, Kubernetes
        nodes, automated updates, SELinux environments.
  \item \textbf{Ubuntu Server}: General-purpose servers, traditional
        deployments, environments requiring package managers.
\end{itemize}
